% titel: Blatt 1, A3: Masse und Bindungsenergie (Fluor-19)
% skript: Mortimer Kap. 2 (Atome, Isotope, Nukleonen)

% Teil I: Werkzeuge, kein Lösungsweg.

Zwischen Massenzahl $A$ und tatsächlicher Atommasse eines monoisotopischen
Elements besteht ein kleiner Massendefekt, der mit der Kernbindungsenergie
zusammenhängt.

\begin{werkzeugkasten}[Nukleonen und Massendefekt]
$A$ = Zahl der Nukleonen (Protonen $Z$ + Neutronen $N$). Relevante Massen
(genauer):
\begin{equation}
  m_p\approx 1.007276\ \mathrm{u},\quad
  m_n\approx 1.008665\ \mathrm{u},\quad
  m_e\approx 0.000549\ \mathrm{u}.
\end{equation}
Beim Kernaufbau wird Energie frei; nach $E=mc^2$ fehlt diese Masse am
Fertigkern. Deshalb ist die Atommasse eines Elements kleiner als die Summe
seiner Einzelteile.

\textit{Skript:} Mortimer Kap. 2.4 (Isotope, Nukleonenmassen), Übung 1/3.
\end{werkzeugkasten}

\begin{achtung}[Darauf achten]
\begin{itemize}
  \item $Z$ und $N$ aus $A$ und Ordnungszahl bestimmen, hier $Z=9$ (Fluor),
        $A=19$ $\Rightarrow N=10$.
  \item Elektronenmasse mitzählen (Atom, nicht Kern) --- klein, aber relevant.
  \item Massendefekt als \emph{relativer} Anteil (%) angeben, nicht nur absolut.
  \item Die häufigste Erklärung in Teil d) ist die freigesetzte
        Bindungsenergie ($\Delta E=\Delta m\,c^2$) --- die anderen Antworten
        sind Ablenker.
\end{itemize}
\end{achtung}
